% !TEX root = ../bb_AiRa_Kappaleet.tex

% ö = \%c3\%b6
% ä = \%C3\%A4

\xdef\sectiontitle{\Isectiontitle} %aiheuttama
\TOCrefs

%%%%%%%%%%%%%%%
\begin{frame}[label=1_11_a]%[label=1_11_TOC1]
\frametitle{\texorpdfstring{\culine[\sectiontitleulinecolor]{\sectiontitle}}{\sectiontitle} }
\label{\TOCname}

\vspace{-0.5cm}

\begin{itemize}[leftmargin=0.5cm,itemsep=2mm]
    \item[1.1]<1-> \hyperlink{1_1_a}{\IaSubsectiontitle}
    \item[1.2]<1-> \hyperlink{1_2_a}{\IbSubsectiontitle}	
    \item[1.3]<1-> \hyperlink{1_3_a}{\IcSubsectiontitle}	
    \item[1.4]<1-> \hyperlink{1_4_a}{\IdSubsectiontitle}
    \item[1.5]<1-> \hyperlink{1_5_a}{\IeSubsectiontitle}
    \item[1.6]<1-> \hyperlink{1_6_a}{\IfSubsectiontitle}
    \item[1.7]<1-> \hyperlink{1_7_a}{\IgSubsectiontitle}
\end{itemize}

\vspace{-1cm}
\begin{itemize}[leftmargin=8.5cm,itemsep=2mm]
    \item[1.8]<1-> \hyperlink{1_8_a}{\IhSubsectiontitle}
	\item[1.9]<1-> \hyperlink{1_9_a}{\IiSubsectiontitle}
	\item[1.10]<1-> \hyperlink{1_10_a}{\IjSubsectiontitle}
	\item[1.11]<1-> \hyperlink{1_11_a}{\CDAlert[blue]{\IkSubsectiontitle}}
\end{itemize}

%\endofslide 
\end{frame}
%%%%%%%%%%%%%%%

\xdef\sectiontitle{\IkSubsectiontitle} 

%%%%%%%%%%%%%%%
\begin{frame}%[label=1_11_a]
\frametitle{\texorpdfstring{\culine[\sectiontitleulinecolor]{\sectiontitle}}{\sectiontitle} }
\framesubtitle{Fononit}
\appear{}

\begin{itemize}[itemsep=1.5mm]
	\item Fononien aallonpituuksille on olemassa alaraja
	\vspace{-1mm}\implies{Käytännössä: \qquad $\frac{1}{2} \appear{\lambda_{\min}}  \equal \appear{a=$} \appear{hilavakio} } 
	\appear{\xdef\loc{south east}%
\begin{tikzpicture}[remember picture,overlay] % anchor=south
    \node (A) [yshift=-0*\figYshift,xshift=\figXshift, anchor=\loc] at (current page.\loc) {\includegraphics[page=1,height=6.7cm]{Figures_booklet/SOM_9_Statistical_mechanics_figures/SOM_Figure_30_cropped.png}};%scale=\figscale. % 8cm max hei
\end{tikzpicture}}
	\item Pienin aallonpituus vastaa suurinta taajuutta
	         \appear{eli suurinta energiaa $E_{\max}$}
\end{itemize}

% 6.5 cm = midway, 10 cm = 3/4 
%\vspace{2.5mm}
\begin{minipage}{8.5cm}
\begin{itemize}[itemsep=1.5mm]
	\item \CDAlert[fuchsia]{Fononien nopeus} ei ole valon nopeus\appear{, vaan \CDAlert[fuchsia]{äänen nopeus $v_{s}$} aineessa}

\item Fononit voivat värähdellä kahdessa \CDAlert[red]{poikittaisessa} suunnassa\appear{, sekä \CDAlert[blue]{pitkittäis-}\\ \CDAlert[blue]{suunnassa}}

\item Olettamalla sekä poikittais- että pitkit\-täis\-fononien nopeudet samoiksi\appear{, fononinen tilatiheys täytyy kertoa tekijällä $3 / 2$}
\end{itemize}
\end{minipage}


\endofslide 
%\endofsection
\end{frame}
%%%%%%%%%%%%%%%

%%%%%%%%%%%%%%%
\begin{frame}
\frametitle{\texorpdfstring{\culine[\sectiontitleulinecolor]{\sectiontitle}}{\sectiontitle} }
\framesubtitle{Fononeihin varastoitunut energia}
\appear{}

\begin{itemize}[itemsep=1.55mm]
	\item Fononien varastoimalle kokonaisenergialle voidaan muodostaa yhtälö:
\[
U_{\text {fononi}}  \equal \appear{\nint_{0}^{E_{\max}}} \appear{E} \appear{\mathbb{N}(E)} \appear{D(E)} \appear{d E}
 \equal 
 \appear{\nint_{0}^{E_{\max }} E} \frac{1}{e^{E / k_{B} T}-1} \frac{12 \pi V}{h^{3} v_{S}{ }^{3}} \appear{E^{2}} \appear{d E}
\]

\item Integraalin laskemiseksi\appear{ meidän täytyy}
% 6.5 cm = midway, 10 cm = 3/4 
%\vspace{2.5mm}
\end{itemize}

\begin{minipage}{7.3cm}
\begin{itemize}[itemsep=1.55mm]
	 \item[] arvioida \Alert[red]{suurin energia $E_{max}$} \appear{ja yhdistää se pienimpään aallonpituuteen \appear{$2 a$}}

\item Tarkastellaan atomeita \Alert[blue]{kuutio\-maisessa kiinteässä kappaleessa} 

\item Atomeita on yhdellä särmällä $N^{1 / 3}$\appear{, ja atomien välinen etäisyys on $L / N^{1 / 3}$} 
\item[] \vspace{-2.2mm}\implies{ Aineessa on fononeita\appear{, joiden pienin aallonpituus on $2 L / N^{1 / 3}$}}

\end{itemize}
\end{minipage}


\xdef\loc{south east}
\begin{tikzpicture}[remember picture,overlay] % anchor=south
    \node (A) [yshift=-0.25*\figYshift,xshift=\figXshift, anchor=\loc] at (current page.\loc) {\includegraphics[page=1,height=5.35cm]{Figures_booklet/SOM_9_Statistical_mechanics_figures/SOM_Figure_31.png}}; %scale=\figscale. % 8cm max hei
\end{tikzpicture}

\endofslide 
%\endofsection
\end{frame}
%%%%%%%%%%%%%%%

%%%%%%%%%%%%%%%%
%\begin{frame}
%\frametitle{\texorpdfstring{\culine[\sectiontitleulinecolor]{\sectiontitle}}{\sectiontitle} }
%
%\begin{itemize}
%	\item Minimum wavelength leads to a maximum wave number $k=\frac{2 \pi}{\lambda}  \equal  \frac{\pi N^{1 / 3}}{L }$ 
%	
%	\item Note, $\mathbf{k}$ is basically a vector (k=$|\mathbf{k}|$)\appear{, it may point at any direction}\appear{, allowed phononic states form a sphere!}
%% 6.5 cm = midway, 10 cm = 3/4 
%\item[] \begin{minipage}{8.5cm}
%\item We consider a cube in $\mathbf{k}$-space, with edge of $\pi N^{1 / 3} / L$, and we set the volume of the cube to be equal to $1 / 8$ of a sphere with radius $k_{\max }$
%\[
%\begin{aligned}
% \bigg( \frac{\pi N^{1 / 3}}{L} \bigg) ^{3} &= \frac{1}{8} \bigg( \frac{4 \pi}{3} k_{\max }^{3} \bigg)  \\
%\Rightarrow \qquad k_{\max } &= \bigg( \frac{3}{4 \pi} 8 \bigg( \frac{\pi N^{1 / 3}}{L} \bigg) ^{3} \bigg) ^{1 / 3} \\
%&= \bigg( \frac{6 \pi^{2} N}{L^{3}} \bigg) ^{1 / 3}= \bigg( \frac{6 \pi^{2} N}{V} \bigg) ^{1 / 3}
%\end{aligned}
%\]
%\end{minipage}
%\vspace{3mm}
%\end{itemize}
%
%
%\xdef\loc{south east}
%\begin{tikzpicture}[remember picture,overlay] % anchor=south
%    \node (A) [yshift=-0.25*\figXshift,xshift=\figXshift, anchor=\loc] at (current page.\loc) {\includegraphics[page=1,height=7cm]{Figures_booklet/SOM_9_Statistical_mechanics_figures/SOM_Figure_32.png}}; %scale=\figscale. % 8cm max hei
%\end{tikzpicture}
%
%
%\endofslide 
%%\endofsection
%\end{frame}
%%%%%%%%%%%%%%%%
%

%%%%%%%%%%%%%%%
\begin{frame}
\frametitle{\texorpdfstring{\culine[\sectiontitleulinecolor]{\sectiontitle}}{\sectiontitle} }
\framesubtitle{Fononien tilatiheys ja energia}
\appear{}

\begin{itemize}[itemsep=1.7mm] \fontsize{11.5}{13}\selectfont %\small
%	\item Minimum wavelength leads to a maximum wave number $k=\frac{2 \pi}{\lambda}  \equal  \frac{\pi N^{1 / 3}}{L }$ 
	\item \CDAlert[blue]{Jaksollisuuden takia} kiinteässä aineessa\appear{ sallitut \Alert[blue]{aaltovektorit $\mathbf{k}$ ovat diskreettejä}:} \appear{yksittäinen fononitila vie} \appear{(ns.~$\mathbf{k}$-avaruudessa}\appear{/tila-avaruudessa)} \appear{tilavuuden  $( 2 \pi /L )^{3}$ }
\end{itemize}
\vspace{1.5mm}
	 
% 6.5 cm = midway, 10 cm = 3/4 
\begin{minipage}{8cm} %\small
\begin{itemize}[itemsep=1.7mm] \fontsize{11.5}{13}\selectfont %\small

\item Huomaa että $\mathbf{k}$ on vektori \appear{(k=$|\mathbf{k}|$)}\appear{, ja voi\\ o\-soit\-taa mielivaltaiseen suuntaan}\appear{, jolloin sallitut fononitilat muodostavat pallon!}

\item \CDAlert[red]{Systeemissä on yhteensä $N$ tilaa}\appear{, jonka täytyy vastata tiheyden} \appear{ja tilavuuden tuloa} \appear{($\mathbf{k}$-avaruudessa):}
\[
\begin{aligned}
N & \equal  \appear{\bigg( \Frac{L}{2\pi} \bigg)^{3}} \appear{\bigg( \Frac{4 \pi}{3} k_{\max }^{3} \bigg)}  \\
\appear{\Rightarrow} \qquad \appear{k_{\max }} & \equal  \appear{\bigg( \frac{8\pi^3}{L^3}} \frac{3N}{4\pi} \appear{\bigg) ^{1 / 3}} \\
& \equal  \appear{\bigg(} \frac{6 \pi^{2} N}{L^{3}} \appear{\bigg) ^{1 / 3}}  \equal  \appear{\bigg(} \frac{6 \pi^{2} N}{V} \appear{\bigg) ^{1 / 3}}
\end{aligned}
\]
\end{itemize}
\end{minipage}


\xdef\loc{south east}
\begin{tikzpicture}[remember picture,overlay] % anchor=south
    \node (A) [yshift=-0.25*\figXshift,xshift=\figXshift, anchor=\loc] at (current page.\loc) {\includegraphics[page=1,height=7cm]{Figures_booklet/SOM_9_Statistical_mechanics_figures/SOM_Figure_32.png}}; %scale=\figscale. % 8cm max hei
\end{tikzpicture}


\endofslide 
%\endofsection
\end{frame}
%%%%%%%%%%%%%%%

%%%%%%%%%%%%%%%
\begin{frame}
\frametitle{\texorpdfstring{\culine[\sectiontitleulinecolor]{\sectiontitle}}{\sectiontitle} }
\framesubtitle{Fononien energia ja Debyen lämpötila}
\appear{}

\begin{itemize}
	\item Nyt kun suurin aaltoluku on johdettu\appear{, voidaan suurin fononin energia arvioida:}
\[
E_{\max}  \equal  \frac{h v_{s} k_{\max }}{2 \pi}  \equal  \frac{h v_{s}}{2 \pi} \appear{\bigg(} \frac{6 \pi^{2} N}{V} \appear{\bigg) ^{1 / 3}}  \equal  \appear{h v_{s}} \appear{\bigg(} \frac{3 N}{4 \pi V} \appear{\bigg) ^{1 / 3}} \appear{\,, \qquad E=\hbar \omega ~\textrm{ja}~\omega = v_s k}
\]

\item Oletimme \appear{että etenevien aaltojen nopeus} \appear{on sama kuin äänen nopeus aineessa}\appear{, mikä on kokeellisesti vahvistettu tosiasia. }\appear{Lisäksi oletimme että \CDAlert[red]{kaikki fononit liikkuvat samalla nopeudella}} 

\item Asettamalla $E_{\max }=k_{B} T_{D}$ \appear{saamme ilmaistua suurimman energian lämpötilana}\appear{, niin sanottuna  \CDAlert[blue]{Debyen lämpötilana}}\appear{, mikä selvästikin riippuu aineen ominaisuuksista}\appear{, kuten $N / V$ ja $v_{s}$:}
\[
T_{D}  \equal \frac{E_{\max }}{k_{B}}  \equal \frac{h v_{S}}{k_{B}} \appear{\bigg(} \frac{3 N}{4 \pi V} \appear{\bigg) ^{1 / 3}}
\]
\end{itemize}

\endofslide 
%\endofsection
\end{frame}
%%%%%%%%%%%%%%%

%%%%%%%%%%%%%%%
\begin{frame}
\frametitle{\texorpdfstring{\culine[\sectiontitleulinecolor]{\sectiontitle}}{\sectiontitle} }
\framesubtitle{Fononeihin varastoitunut energia ja Debyen funktio}

\begin{itemize}
	\item Nyt kun maksimienergia on laskettu\appear{, voimme palata alkuperäiseen kokonaisenergian integraaliin:}
	\appear{$$
	U_{\text {fononi}}  \equal   \appear{\nint_{0}^{k_{B} T_{D}}} \appear{E} \fracc{1}{e^{E / k_{B} T}-1} \fracc{9 N}{ \big( k_{B} T_{D} \big) ^{3}} \appear{E^{2}} \appear{d E}  \equal \appear{9 k_{B} N} \frac{T^{4}}{T_{D}^{3}} \appear{\nint_{0}^{T_{D} / T}} \frac{x^{3}}{e^{x}-1} \appear{d x}
	$$
	\appear{missä muuttuja $E$ on vaihdettu} \appear{$x$:ksi} $\appear{(x} \equiv \appear{E/k_{B} T)}$}
\item Moolille ainetta $N=N_{A}$\appear{, ja koska $k_{B} N_{A}=R$}\appear{, saamme lopuksi:}
\[
U_{\text{fononi,\,per mooli}}  \equal \appear{9 R} \frac{T^{4}}{T_{D}^{3}} \appear{\nint_{0}^{T_{D} / T}} \frac{x^{3}}{e^{x}-1} \appear{d x} \qquad
\]

\item Valitettavasti \CDAlert[red]{Debyen funktion} \Alert[tunired]{integraalia ei voida analyyttisesti laskea}\appear{, vaan se täytyy arvioida numeerisesti}\appear{, ja verrata kokeellisiin tuloksiin}
\end{itemize}


\endofslide 
%\endofsection
\end{frame}
%%%%%%%%%%%%%%%

%%%%%%%%%%%%%%%
\begin{frame}
\frametitle{\texorpdfstring{\culine[\sectiontitleulinecolor]{\sectiontitle}}{\sectiontitle} }
\framesubtitle{Fononien vaikutus ominaislämpökapasiteetteihin}
\appear{}

\begin{itemize}
	\item Ylläkuvattu \CDAlert[fuchsia]{Debyen malli} \appear{\Alert[fuchsia]{täsmää hyvin kokeellisten lämpökapasiteettien kanssa}}\appear{, kuten kuvasta nähdään}
	
	\item \appear{\CDAlert[red]{Johtaminen oli nopea ja yksinkertaistettu}}\appear{, parempi johtaminen suoritetaan johdantokurssilla kiinteän olomuodon fysiikkaan!}
\end{itemize}

\xdef\loc{south}
\begin{tikzpicture}[remember picture,overlay] % anchor=south
    \node (A) [yshift=-0.25*\figYshift,xshift=\figXshift, anchor=\loc] at (current page.\loc) {\includegraphics[page=1,width=14cm]{Figures_booklet/SOM_9_Statistical_mechanics_figures/SOM_Figure_33.png}}; %scale=\figscale. % 8cm max hei
\end{tikzpicture}


%\endofslide 
\endofsection
\end{frame}
%%%%%%%%%%%%%%%
